\chapter{Phase D --- Standard examples (early subset)}

Phase~D, in full, will populate the standard zoo of affine algebraic groups
($\Ga$, $\Gm$, $\mu_n$, $\alpha_p$, $\GL_n$, $\SL_n$, $T_n$, $B_n$, $U_n$).
The early subset implemented in parallel with Phase~B comprises $\Ga$, $\Gm$,
and $\GL_n$: the first two depend only on the Phase~A typeclass framework,
and $\GL_n$ is required as the target of the embedding theorem (Phase~B4).

\section{The additive group $\Ga$}

\begin{definition}[Additive group scheme $\Ga$]\label{def:Ga}\issue{6}
\uses{def:algebraic-group}
For a scheme $S$, the \emph{additive group scheme over $S$} is
\[
  \Ga{}_{,S} \;=\; S \times_{\Spec \mathbb{Z}} \Spec \mathbb{Z}[X],
\]
viewed as an affine group scheme over $S$ via the additive Hopf structure on
the underlying polynomial ring:
\[
  \Delta(X) = X \otimes 1 + 1 \otimes X, \qquad
  \varepsilon(X) = 0, \qquad
  S(X) = -X.
\]
This is an algebraic group over $S$.
\end{definition}

\begin{lemma}[$\Ga(R) = R$]\label{lem:Ga-points}\issue{6}
\uses{def:Ga, def:points}
For a field $k$ and a $k$-algebra $R$, there is a natural isomorphism of
additive groups $\Ga{}_{,\Spec k}(R) \cong (R, +)$.
\end{lemma}

\section{The multiplicative group $\Gm$}

\begin{definition}[Multiplicative group scheme $\Gm$]\label{def:Gm}\issue{7}
\uses{def:algebraic-group}
For a scheme $S$, the \emph{multiplicative group scheme over $S$} is
\[
  \Gm{}_{,S} \;=\; S \times_{\Spec \mathbb{Z}} \Spec \mathbb{Z}[X, X^{-1}],
\]
viewed as an affine group scheme over $S$ via the multiplicative Hopf
structure on Laurent polynomials:
\[
  \Delta(X) = X \otimes X, \qquad
  \varepsilon(X) = 1, \qquad
  S(X) = X^{-1}.
\]
This is an algebraic group over $S$.
\end{definition}

\begin{lemma}[$\Gm(R) = R^\times$]\label{lem:Gm-points}\issue{7}
\uses{def:Gm, def:points}
For a field $k$ and a $k$-algebra $R$, there is a natural isomorphism of
commutative groups $\Gm{}_{,\Spec k}(R) \cong R^\times$.
\end{lemma}

\section{The general linear group $\GL_n$}

\begin{definition}[General linear group scheme $\GL_n$]\label{def:GLn}\issue{8}
\uses{def:algebraic-group}
For a positive integer $n$ and a scheme $S$, the \emph{general linear group
scheme of rank~$n$ over $S$} is
\[
  \GL_n{}_{,S} \;=\; S \times_{\Spec \mathbb{Z}} \Spec \mathbb{Z}[x_{ij}, t]/(t \cdot \det(x_{ij}) - 1),
\]
viewed as an affine group scheme over $S$ via the matrix-multiplication Hopf
structure:
\[
  \Delta(x_{ij}) = \sum_{k=1}^n x_{ik} \otimes x_{kj}, \qquad
  \varepsilon(x_{ij}) = \delta_{ij}, \qquad
  S(x_{ij}) = \det(x_{ij})^{-1} \cdot (\text{cofactor matrix})_{ji}.
\]
This is an algebraic group over $S$.
\end{definition}

\begin{lemma}[$\GL_n(R) = \mathrm{GL}_n(R)$ as abstract group]\label{lem:GLn-points}\issue{8}
\uses{def:GLn, def:points}
For a field $k$ and a $k$-algebra $R$, there is a natural isomorphism of
groups $\GL_n{}_{,\Spec k}(R) \cong \mathrm{GL}_n(R)$, where the right side
is the group of invertible $n \times n$ matrices over $R$
(Mathlib's \texttt{Matrix.GeneralLinearGroup (Fin n) R}).
\end{lemma}

\begin{remark}\label{rem:GLn-as-embedding-target}
$\GL_n$ is the target of the embedding theorem stated in
Theorem~\ref{thm:embedding-statement}. Phase~B4 will promote that
statement to a theorem; once that is done, Phase~B5 will package the
corollary that every algebraic group is a closed subgroup of some
$\GL_n$.
\end{remark}
